% -*- coding: utf-8 -*-
% !TEX program = xelatex
\documentclass[plain,noanswer,medmath]{../../template/jnuexam}
\usepackage{../../template/kysx}

\begin{document}


\examtitle{2011}{1}


\makepart{选择题}{1～8小题，每小题4分，共32分}


\begin{problem}
曲线$y = (x - 1)(x - 2)^2(x - 3)^3(x - 4)^4$的拐点是\pickout{}
\begin{abcd}
\item $(1,0)$．
\item $(2,0)$．
\item $(3,0)$．
\item $(4,0)$．
\end{abcd}
\end{problem}


\begin{problem}
设数列$\left\{a_n \right\}$单调减少，$\lim_{n \to \infty} a_n = 0$，
$S_n = \sum_{k = 1}^n a_k \;(n = 1,2, \cdots)$无界，
则幂级数$\sum_{n = 1}^{\infty}a_n (x - 1)^n$的收敛域为\pickout{}
\begin{abcd}
\item $(-1,1]$．
\item $[-1,1)$．
\item $[0,2)$．
\item $(0,2]$．
\end{abcd}
\end{problem}


\begin{problem}
设函数$f(x)$具有二阶连续导数，且$f(x) > 0$，$f'(0) = 0$，
则函数$z = f(x)\ln f(y)$在点$(0,0)$处取得极小值的一个充分条件是\pickout{}
\begin{abcd}
\item $f(0) > 1$，$f''(0) > 0$．
\item $f(0) > 1$，$f''(0) < 0$．
\item $f(0) < 1$，$f''(0) > 0$．
\item $f(0) < 1$，$f''(0) < 0$．
\end{abcd}
\end{problem}


\begin{problem}
设$I = \int_{0}^{\frac{\pi}{4}}\ln\sin x\dx$，$J = \int_{0}^{\frac{\pi}{4}}\ln\cot x\dx$，
$K = \int_{0}^{\frac{\pi}{4}}\ln\cos x\dx$，则$I,\;J,\;K$的大小关系是\pickout{}
\begin{abcd}
\item $I < J < K$．
\item $I < K < J$．
\item $J < I < K$．
\item $K < J < I$．
\end{abcd}
\end{problem}


\begin{problem}
设$A$为$3$阶矩阵，将$A$的第$2$列加到第$1$列得矩阵$B$，
再交换$B$的第$2$行与第$3$行得单位矩阵，记$P_1 = \begin{pmatrix}
  1 & 0 & 0 \\
  1 & 1 & 0 \\
  0 & 0 & 1
\end{pmatrix}$，$P_2 = \begin{pmatrix}
  1 & 0 & 0 \\
  0 & 0 & 1 \\
  0 & 1 & 0
\end{pmatrix}$，则$A =$\pickout{}
\begin{abcd}
\item $P_1P_2$．
\item $P_1^{-1}P_2$．
\item $P_2P_1$．
\item $P_2 P_1^{-1}$．
\end{abcd}
\end{problem}


\begin{problem}
设$A = (\alpha_1,\alpha_2,\alpha_3,\alpha_4)$是$4$阶矩阵，$A^*$为$A$的伴随矩阵，
若$(1,0,1,0)^T$是方程组$Ax = 0$的一个基础解系，则$A^* x = 0$的基础解系可为\pickout{}
\begin{abcd}
\item $\alpha_1,\alpha_3$．
\item $\alpha_1,\alpha_2$．
\item $\alpha_1,\alpha_2,\alpha_3$．
\item $\alpha_2,\alpha_3,\alpha_4$．
\end{abcd}
\end{problem}


\begin{problem}
设$F_1(x)$，$F_2(x)$为两个分布函数，其相应的概率密度$f_1(x)$，$f_2(x)$是连续函数，
则必为概率密度的是\pickout{}
\begin{abcd}
\item $f_1(x) f_2(x)$．
\item $2f_2(x) F_1(x)$．
\item $f_1(x) F_2(x)$．
\item $f_1(x) F_2(x) + f_2(x) F_1(x)$．
\end{abcd}
\end{problem}


\begin{problem}
设随机变量$X$与$Y$相互独立，且$E(X)$与$E(Y)$存在，
记$U = \max\{X,Y\}$，$V = \min\{X,Y\}$，则$E(UV) =$\pickout{}
\begin{abcd}
\item $E(U) \cdot E(V)$．
\item $E(X) \cdot E(Y)$．
\item $E(U) \cdot E(Y)$．
\item $E(X) \cdot E(V)$．
\end{abcd}
\end{problem}


\makepart{填空题}{9～14小题，每小题4分，共24分}


\begin{problem}
曲线$y = \int_0^x \tan t\dt$（$0 \le x \le \frac{\pi}{4}$）的弧长$s =$ \fillin{}．
\end{problem}


\begin{problem}
微分方程$y' + y = \e^{- x}\cos x$满足条件$y(0) = 0$的解为$y =$ \fillin{}．
\end{problem}


\begin{problem}
设函数$F(x,y) = \int_0^{xy} \frac{\sin t}{1 + t^2}\dt$，
则$\left. \frac{\pd^2 F}{\pd x^2} \right|_{\begin{subarray}{l}
  x = 0 \\
  y = 2
\end{subarray}} =$ \fillin{}．
\end{problem}


\begin{problem}
设$L$是柱面方程$x^2 + y^2 = 1$与平面$z = x + y$的交线，从$z$轴正向往$z$轴负向看去为逆时针方向，
则曲线积分$\oint_L xz\dx + x\dy + \frac{y^2}{2}\dz =$ \fillin{}．
\end{problem}


\begin{problem}
若二次曲面的方程为$x^2 + 3y^2 + z^2 + 2axy + 2xz + 2yz = 4$，
经过正交变换化为$y_1^2 + 4z_1^2 = 4$，则$a =$ \fillin{}．
\end{problem}


\begin{problem}
设二维随机变量$(X,Y)$服从正态分布$N\left(\mu,\mu;\sigma^2,\sigma^2;0\right)$，
则$E\left(XY^2\right) =$ \fillin{}．
\end{problem}


\makepart{解答题}{15～23小题，共94分}


\begin{problem}[points=10]
求极限$\lim_{x \to 0} \left(\frac{\ln(1 + x)}{x} \right)^{\frac{1}{\e^x - 1}}$．
\end{problem}


\begin{problem}[points=9]
设函数$z = f(xy,yg(x))$，其中函数$f$具有二阶连续偏导数，
函数$g(x)$可导且在$x = 1$处取得极值$g(1) = 1$，
求$\left. \frac{\pd^2 z}{\pdx\pd y} \right|_{\begin{subarray}{l}
  x = 1 \\
  y = 1
\end{subarray}}$．
\end{problem}


\begin{problem}[points=10]
求方程$k\arctan x - x = 0$不同实根的个数，其中$k$为参数．
\end{problem}


\begin{problem}[points=10]
\begin{enumerate}
\item
证明：对任意的正整数$n$，都有$\frac{1}{n + 1} < \ln \left(1 + \frac{1}{n} \right) < \frac{1}{n}$成立．
\item
设$a_n = 1 + \frac{1}{2} + \cdots + \frac{1}{n} - \ln n$（$n = 1,2, \cdots$），
证明数列$\left\{a_n \right\}$收敛．
\end{enumerate}
\end{problem}


\begin{problem}[points=11]
已知函数$f(x,y)$具有二阶连续偏导数，且
$$f(1,y) = 0,\quad f(x,1) = 0,\quad \iint_D f(x,y)\dx\dy = a,$$
其中$D = \left\{(x,y)\middle|0 \le x \le 1,0 \le y \le 1 \right\}$，
计算二重积分$I = \iint_D xy f''_{xy}(x,y)\dx\dy$．
\end{problem}


\begin{problem}[points=11]
设向量组$\alpha_1 = (1,0,1)^T$, $\alpha_2 = (0,1,1)^T$, $\alpha_3 = (1,3,5)^T$
不能由向量组$\beta_1 = (1,1,1)^T$, $\beta_2 = (1,2,3)^T$, $\beta_3 = (3,4,a)^T$线性表示．\par
\begin{enumerate*}
\item 求$a$的值；
\item 将$\beta_1,\beta_2,\beta_3$用$\alpha_1,\alpha_2,\alpha_3$线性表示．
\end{enumerate*}
\end{problem}


\begin{problem}[points=11]
$A$为三阶实对称矩阵，$A$的秩$r(A) = 2$，且$A\begin{pmatrix}
   1 & 1 \\
   0 & 0 \\
  -1 & 1
\end{pmatrix} = \begin{pmatrix}
  -1 & 1 \\
   0 & 0 \\
   1 & 1
\end{pmatrix}$．\par
\begin{enumerate*}
\item 求$A$的特征值与特征向量；
\item 求矩阵$A$．
\end{enumerate*}
\end{problem}


\begin{problem}[points=11]
设随机变量$X$与$Y$的概率分布分别为
$$\begin{array}{|c|5cc|}
\hline
  X &           0 & 1 \\
\hline
  P & \frac{1}{3} & \frac{2}{3} \\
\hline
\end{array}\quad \quad \begin{array}{|c|5ccc|}
\hline
  Y &          -1 &           0 & 1 \\
\hline
  P & \frac{1}{3} & \frac{1}{3} & \frac{1}{3} \\
\hline
\end{array}$$
且$P\left\{X^2 = Y^2 \right\} = 1$．
\begin{enumerate}
\item 求二维随机变量$(X,Y)$的概率分布；
\item 求$Z = XY$的概率分布；
\item 求$X$与$Y$的相关系数$\rho_{XY}$．
\end{enumerate}
\end{problem}


\begin{problem}[points=11]
设$X_1$, $X_2$, $\cdots$, $X_n$为来自正态总体$N(\mu_0,\sigma^2)$的简单随机样本，
其中$\mu_0$已知，$\sigma^2 > 0$未知．$\overline{X}$和$S^2$分别表示样本均值和样本方差．\par
\begin{enumerate*}
\item 求参数$\sigma^2$的最大似然估计量$\hat{\sigma}^2$；
\item 计算$E(\hat{\sigma}^2)$和$D(\hat{\sigma}^2)$．
\end{enumerate*}
\end{problem}


\end{document}
